\documentclass[twocolumn]{aastex631}
% Auto-generated by jansky_research.vlasspm._write_macros -- do not edit.
% vpmReal* come from results/vlasspm_metrics.json + vlasspm_vetting.json (real leg);
% vpmSyn* from the offline fixture. The inactive namespace is written as placeholders
% and preserve_live_macros keeps the other leg's live values.
\newcommand{\vpmSource}{real: VLASS Quick-Look epochs 1, 2, 3 (QL3.1+3.2) and 4.1 component catalogues}
\newcommand{\vpmSynNMovers}{25}
\newcommand{\vpmSynNRecovered}{17}
\newcommand{\vpmSynNFalse}{0}
\newcommand{\vpmSynCompleteness}{0.50}
\newcommand{\vpmSynNullCand}{0.00}
\newcommand{\vpmSynKTrue}{0.10}
\newcommand{\vpmSynKFit}{0.10}
\newcommand{\vpmSynTailOld}{0.41}
\newcommand{\vpmSynTailFit}{0.009}
\newcommand{\vpmRealNEOne}{1{,}874{,}639}
\newcommand{\vpmRealNETwo}{1{,}873{,}519}
\newcommand{\vpmRealNEThree}{2{,}360{,}730}
\newcommand{\vpmRealNEFour}{1{,}081{,}153}
\newcommand{\vpmRealTStart}{2017.7}
\newcommand{\vpmRealTEnd}{2026.1}
\newcommand{\vpmRealAreaDeg}{33{,}838}
\newcommand{\vpmRealFloorMin}{0.08}
\newcommand{\vpmRealFloorMax}{0.31}
\newcommand{\vpmRealNStaticPairs}{4.9}
\newcommand{\vpmRealKStruct}{0.10}
\newcommand{\vpmRealTailLegacy}{0.041}
\newcommand{\vpmRealTailMeas}{0.032}
\newcommand{\vpmRealTailFit}{0.0103}
\newcommand{\vpmRealTailFitFive}{0.0012}
\newcommand{\vpmRealMedZFit}{0.81}
\newcommand{\vpmRealTailSouth}{0.025}
\newcommand{\vpmRealTailSouthRatio}{2.3}
\newcommand{\vpmRealUVMu}{3.45}
\newcommand{\vpmRealUVGaiaMu}{3.23}
\newcommand{\vpmRealUVResid}{1.79}
\newcommand{\vpmRealUVCompA}{0.41}
\newcommand{\vpmRealUVCompB}{0.00}
\newcommand{\vpmRealUVCompC}{1.46}
\newcommand{\vpmRealNCandBefore}{11}
\newcommand{\vpmRealNCandAfter}{1}
\newcommand{\vpmRealNStaticVetted}{10}
\newcommand{\vpmRealNullCand}{0.02}
\newcommand{\vpmRealNStars}{303}
\newcommand{\vpmRealNStarChance}{16}
\newcommand{\vpmRealCutThr}{0.9}
\newcommand{\vpmRealCutThrAll}{1.6}
\newcommand{\vpmRealNInjCal}{40{,}000}
\newcommand{\vpmRealNRecCal}{23{,}932}
\newcommand{\vpmRealKeepAll}{0.984}
\newcommand{\vpmRealKeepLowSNR}{0.951}
\newcommand{\vpmRealKeepSlow}{0.93}
\newcommand{\vpmRealStaticRej}{14}
\newcommand{\vpmRealStaticRejAll}{6.6}
\newcommand{\vpmRealNInj}{20{,}000}
\newcommand{\vpmRealCompThreeCut}{0.964}
\newcommand{\vpmRealCompThreeNoCut}{0.965}
\newcommand{\vpmRealCompOneFiveCut}{0.952}
\newcommand{\vpmRealCompOneFiveNoCut}{0.961}
\newcommand{\vpmRealCompSlowThreeCut}{0.08}
\newcommand{\vpmRealCompFloorThreeCut}{0.00}
\newcommand{\vpmRealRateLo}{0.53}
\newcommand{\vpmRealRateMid}{0.92}
\newcommand{\vpmRealRateMin}{0.3}
\newcommand{\vpmRealRateMax}{5}
\newcommand{\vpmRealLimThree}{9.2\times10^{-5}}
\newcommand{\vpmRealLimOneFive}{9.3\times10^{-5}}
\newcommand{\vpmRealAllSky}{3.8}
\newcommand{\vpmRealNVetted}{11}
\newcommand{\vpmRealNNewMovers}{0}
\newcommand{\vpmRealNSouthVetted}{5}

\newcommand{\code}[1]{\texttt{#1}}
\shorttitle{A blind VLASS proper-motion search}
\shortauthors{Barbere}
\begin{document}
\title{No Optically Dark Compact Radio Movers in Four VLASS Epochs}
\author[0009-0008-3289-4447]{Joseph Barbere}
\affiliation{Independent researcher}

\begin{abstract}
We searched three VLASS Quick Look epochs over the whole survey, and a fourth over
\vpmRealAreaFour~deg$^2$, for moving radio sources, consulting optical and infrared catalogues only
at the end. The search recovers the unresolved UV/BL~Ceti binary at
$\vpmRealUVFitMu\pm\vpmRealUVFitErr$~arcsec~yr$^{-1}$ (system motion \vpmRealSysMu~arcsec~yr$^{-1}$)
and finds no other mover. With errors calibrated on static sources, a compactness cut calibrated
on nearby Gaia stars leaves UV~Ceti the only one of \vpmRealNCandBefore\ candidates; the rest are
static in the images. For persistent (three-epoch) compact sources of
\vpmRealFluxLoThree--\vpmRealFluxHiThree~mJy beyond \vpmRealPlxDomain~pc, moving at
\vpmRealRateDomLo--\vpmRealRateMax~arcsec~yr$^{-1}$ with no co-moving Gaia or CatWISE counterpart,
the 95\% upper limit is $\vpmRealLimDomThree$~deg$^{-2}$.
\end{abstract}

\keywords{Proper motions (1295) --- Radio astrometry (1337) --- Sky surveys (1464) ---
Radio continuum emission (1340)}

\section{Search}
\label{sec:search}
Radio proper motions from survey data have so far been measured for stars selected by their Gaia
proper motions \citep{driessen2023,de2024}; \citet{atri2022} measured VLBA proper motions of
compact, flat-spectrum, variable Galactic-plane sources, a targeted rather than a blind search. A
blind search could find movers invisible to Gaia, such as nearby cool brown dwarfs. The VLA Sky
Survey \citep[VLASS;][]{lacy2020} has covered the sky north of Dec $-40^\circ$ three times at
3~GHz with a 2.5-arcsec beam, and a fourth epoch covers \vpmRealAreaFour~deg$^2$ so far. We used
the CIRADA component tables for epochs 1 and 2 \citep{gordon2021} and the NRAO Quick Look lists
for epochs 3 and 4, all fitted by PyBDSF \citep{pybdsf}: after the quality cuts they hold
\vpmRealNEOne, \vpmRealNETwo, \vpmRealNEThree\ and \vpmRealNEFour\ components observed between
\vpmRealTStart\ and \vpmRealTEnd.

A component with no counterpart within 2.5~arcsec in another epoch is an orphan, and only
orphans with no other component within 30~arcsec in their own epoch are used, because split
extended sources produce correlated orphans. Orphans in two epochs are linked when the implied
rate lies between \vpmRealRateMin\ and \vpmRealRateMax~arcsec~yr$^{-1}$ and the shift is
significant at $3\sigma$; a third epoch must then contain an orphan within $3\sigma$ of the
extrapolated straight line at its own observing date. Every triple of epochs is searched.

Each component's positional error is the elliptical Gaussian error of \citet{condon1997}
\citep[see also][]{condon1998,prandoni2000}, from its fitted shape, beam and peak
signal-to-noise. To it we add an astrometric floor
measured from bright static sources in declination bands
(\vpmRealFloorMin--\vpmRealFloorMax~arcsec) and a structure term: a resolved source's centroid
moves between epochs by a fraction \vpmRealKStruct\ of its deconvolved size along each of its
axes. That fraction was fitted to \vpmRealNStaticPairs~million pairs of isolated static sources
so that their $3\sigma$ tail matches a Rayleigh distribution; on a held-out half of the sky the
fraction beyond $3\sigma$ falls from \vpmRealTailMeas\ with the catalogue errors to
\vpmRealTailFit\ (expected 0.011), though it remains \vpmRealTailSouthRatio\ times too large south
of Dec $-35^\circ$.

An earlier version of the search returned almost only resolved static sources, which led us to
add a compactness cut (deconvolved over beam major axis), with threshold and rule fixed on
injections before it was applied. Faint point sources are not fitted with zero size, so injected
sizes were drawn, by signal-to-noise, from \vpmRealNStars\
detections of stars within 100~pc in the Gaia Catalogue of Nearby Stars \citep{gcns2021}. The criterion was the smallest threshold keeping at least 95\% of the recovered injections in
every signal-to-noise bin, and the rule that rejects more static sources. Two of three detections below \vpmRealCutThr\ keep \vpmRealKeepAll\ of the
injections (\vpmRealKeepLowSNR\ below signal-to-noise 7) and remove \vpmRealStaticRej\% of static
sources, more than requiring all three does.

Rotating the later epochs by a few arcminutes in right ascension measures chance alignments of
unrelated orphans; it cannot see a static source whose fitted centroid shifts between epochs.
\vpmRealNInj\ movers injected at random sky positions with the same errors, size noise and annual
parallax give the completeness.

\section{Results}
\label{sec:results}
The mover found in epochs 2--4 at the position of UV~Ceti is the unresolved Luyten 726-8 system:
UV and BL~Cet are \vpmRealCompSep~arcsec apart and both have Gaia RUWE of
\vpmRealRuweMin--\vpmRealRuweMax. A straight line through the three detections, weighted by their
covariances and with the Gaia DR3 \citep{gaiadr3} parallax fixed, gives
$\vpmRealUVFitMu\pm\vpmRealUVFitErr$~arcsec~yr$^{-1}$. It differs from the system motion
(\vpmRealSysMu~arcsec~yr$^{-1}$; UCAC4, \citealt{ucac4}, via SIMBAD, \citealt{simbad}) by a vector
of \vpmRealSysOffset~arcsec~yr$^{-1}$ ($\chi^2=\vpmRealSysChi$ for 2 degrees of freedom), comparable
with the \vpmRealOrbScale~arcsec~yr$^{-1}$ between the two components' Gaia motions, so orbital
motion and blending of the two stars account for it. Two of its three compactness values
(\vpmRealUVCompA, \vpmRealUVCompB, \vpmRealUVCompC) pass the cut.

Before the compactness cut the search returns \vpmRealNCandBefore\ candidates. Cutouts of the
other \vpmRealNStaticVetted\ show emission persisting at the first-epoch position, so they are
static. The cut removes all \vpmRealNStaticVetted, two narrowly (\vpmRealMarginAName\
and \vpmRealMarginBName\ would pass at \vpmRealMarginAThr\ and \vpmRealMarginBThr), and because the
pre-stated two-of-three rule won on the static population, no by-eye step lies between the
catalogues and the limit. A candidate is optically dark if no Gaia DR3 or CatWISE2020
\citep{catwise2020} source lies within $3\sigma$ of the track's position at the catalogue epoch
(with a 0.5-arcsec floor added in quadrature) or moves within 30\% of the radio rate. Only UV~Ceti
reaches this step, so chance coincidence cannot have removed a mover.

Completeness is near zero below \vpmRealRateLo~arcsec~yr$^{-1}$, where a mover stays within the
static match radius, climbs steeply between \vpmRealFineEdgeLo\ and
\vpmRealFineEdgeHi~arcsec~yr$^{-1}$, and is at least \vpmRealCompPlxNone\ in every bin from
\vpmRealRateDomLo\ to \vpmRealRateMax~arcsec~yr$^{-1}$ (Figure~\ref{fig:comp}). Parallax lowers the
worst of those bins to \vpmRealCompPlxFour, \vpmRealCompPlxEight\ and \vpmRealCompPlxSixteen\ at
4, 8 and 16~pc; we adopt \vpmRealPlxDomain~pc, beyond which the loss is below 5\%, meaning
tangential velocities above \vpmRealVtanMin~km~s$^{-1}$. With no new mover and no rate prior, the
worst bin gives 95\% upper limits over the \vpmRealAreaDeg~deg$^2$ searched in three epochs of
$\vpmRealLimDomThree$~deg$^{-2}$ for sources of \vpmRealFluxLoThree--\vpmRealFluxHiThree~mJy per
epoch beyond \vpmRealPlxDomain~pc, and $\vpmRealLimDomOneFive$~deg$^{-2}$ for
\vpmRealFluxLoOneFive--\vpmRealFluxHiOneFive~mJy without parallax, or about \vpmRealAllSky\ over the
whole sky if movers are isotropic. The heavy error tails lower the true completeness by an
estimated \vpmRealTailBoundThree\% at \vpmRealFluxThree~mJy and \vpmRealTailBoundOneFive\% at
\vpmRealFluxOneFive~mJy. The limit applies only to movers unresolved at 2.5~arcsec: the cut
removes a mover with resolved emission and a binary a few tenths of a beam wide fitted as one
elongated source in two epochs.

\begin{figure}
\plotone{figures/vlasspm_completeness.pdf}
\caption{Fraction of injected point-source movers recovered in each rate bin, at \vpmRealFluxThree\
and \vpmRealFluxOneFive~mJy with (solid) and without (dashed) the compactness cut, without parallax.
Points with horizontal bars are finer bins across the completeness edge
(\vpmRealFluxThree~mJy, with the cut). The dotted line marks the three-epoch rate of UV~Ceti.}
\label{fig:comp}
\end{figure}

Reproducibility: \code{jansky\_research.vlasspm} and \code{scripts/vlasspm\_*.py}; every number
derived from the data is pipeline-generated (\code{results/vlasspm\_*.json} $\to$ macros). The
VLASS catalogues and cutouts, Gaia DR3, the Gaia Catalogue of Nearby Stars and CatWISE2020 are
public and need no authentication.

\section*{AI use disclosure}
This paper was developed collaboratively by the author with Anthropic's Claude
\citep{anthropicclaude} (models Claude Opus~4.x and 5.5, accessed 2026). Claude contributed to
the design, implementation, and testing of the \code{vlasspm} module and the drafting of this
manuscript. The author directed the work, independently reviewed and verified all code, results,
and citations, and takes full responsibility for the content. An AI/LLM is not an author.

\begin{acknowledgments}
We used the VLASS Quick Look images and catalogues (NRAO, with CIRADA and CADC), Gaia DR3 and
the Gaia Catalogue of Nearby Stars (ESA), and CatWISE2020; the analysis builds on the
\code{jansky} project. The author acknowledges the collaborative contribution of Anthropic's
Claude.
\end{acknowledgments}

\software{\code{jansky-research} \citep{janskyresearch}, Claude \citep{anthropicclaude},
PyBDSF \citep{pybdsf}, astroquery, Astropy, NumPy, SciPy, Matplotlib, jansky \citep{jansky}}

\bibliography{refs}
\end{document}
